\documentclass[11pt]{jsarticle}

%パッケージのinclude
\usepackage{amsmath} % 幅広い数式を提供
\usepackage{txfonts} % 数式のTimes系フォントを実現するため
\usepackage{listings, jlisting} % プログラムのソースコードを書くため
\usepackage{ascmac} % 枠
\usepackage{fancybox}
\usepackage[dvipdfmx]{graphicx}
\usepackage{float} % 図挿入時に[H]を使うため
\usepackage{array}

\renewcommand{\lstlistingname}{ソースコード}
\lstset{ %ソースコード貼り付けのために
basicstyle={\ttfamily},
identifierstyle={\small},
commentstyle={\smallitshape},
keywordstyle={\small\bfseries},
ndkeywordstyle={\small},
stringstyle={\small\ttfamily},
frame={tb},
breaklines=true,
columns=[l]{fullflexible},
numbers=left,
xrightmargin=0zw,
xleftmargin=3zw,
numberstyle={\scriptsize},
stepnumber=1,
numbersep=1zw,
lineskip=-0.5ex
} % jlisting.styファイル設置を忘れずに...

% タイトル情報
\title{シミュレータ仕様書}
\author{高橋 優作}
\date{\today}

\begin{document}

\begin{titlepage}
    \centering
    \vspace*{35mm}
    {\Huge \bfseries シミュレータ仕様書 \par}
    \vspace{25mm}
    {\Large 高橋 優作 \par}
    \vspace{10mm}
    {\Large \today \par}
    \vfill
\end{titlepage}

\tableofcontents
\newpage

\section{はじめに}
こんにちは．

\section{クラス図}


\section{詳細仕様}
\subsection{グローバル変数設定}

\begin{lstlisting}[caption=Global Variable]
# データ設定
U = list(range(200))    # クライアント生成
D = list(range(5))      # UAV生成
I = list(range(49))     # 領域生成
M = 5                   # 一度に通信できるクライアント数
width = 7               # 横
height = 7              # 縦
move_k = 2              # UAVはmove_kマス移動可能
MOVE_p = 0.3            # クライアントはMOVE_pの確率で移動する
T_MAX = 200             # 自己最適モードの最大時刻
B = 24                  # サーバ設定

# パラメータ
alpha = 0.1             # 移動距離ペナルティ
beta = 0.3              # 社会最適における重複配置ペナルティ
target_bias = 0.3       # 自己最適で近傍にUEがいないとき，未収集UEへ近づくための係数
\end{lstlisting}

\subsection{Areaクラス}
小領域を管理するクラスである．本クラスのインスタンス変数は以下の通りである．
\begin{itemize}
    \item \verb|id| : 領域番号
    \item \verb|clients| : 領域内のクライアントを格納する集合
    \item \verb|neighbors| : 隣接領域を格納するリスト
\end{itemize}

本クラスのメソッドは以下の通りである．
\subsubsection{\texttt{add\_client(u)}}
クライアント \(u\) を当該領域に追加するメソッドである．
Fieldクラス内の \verb|set_clients| メソッドで使用されている．

\begin{lstlisting}[caption=Areaクラス]
class Area:
    def __init__(self, area_id):
        self.id = area_id
        self.clients = set()
        self.neighbors = []

    def add_client(self, u):
        self.clients.add(u)
\end{lstlisting}

\subsection{Fieldクラス}
Fieldクラスは各小領域の隣接領域の取得や，クライアントが動いた際のリアルタイムなクライアント情報を管理する．

\subsection{UAVクラス}


\subsection{Serverクラス}


\subsection{Gameクラス}



\end{document}